\begin{english}
\vspace*{\fill} % Vertically center the page content
\section*{Abstract}
\label{AbstractEnglish}

This thesis aims to solve the problem of human speech recognition using visual cues, more specifically a video of the speech in question. This is a problem very similar to the classic problem of Automatic Speech Recognition (ASR) and is thereby often called Visual Speech Recognition (VSR) or Audio-Visual Speech Recognition (AVSR), in the cases where both audio as well as video cues are used to predict speech. The field of ``lip reading'' has grown immensely in the last few years, especially since the advent of Transformers and the ability to train large artificial intelligence models using deep learning. Recent studies tend to divide the problem of VSR into two separate sub-problems; first, the prediction of phonemes (the smallest units containing phonetic information for human speech) and second, the combination of the predicted phonemes so as to predict words and, ultimately, phrases that the original speaker uttered.

During this particular work, more emphasis is placed on solving the first sub-problem, more specifically the problem of phoneme prediction using the video of a speaker. An artificial intelligence model is developed, which takes in as input the video of the speaker, centered around the area of the lips and without auditory information and returns its predicted phonemes per video frame. This model contains a 3D Convolutional Neural Network (CNN) responsible for the extraction of visual features, as well as a BiLSTM architecture as a Sequence Encoder. Moreover, a similar model is implemented which takes in as input both the video of the area of the lips as well as audio cues. This audio-visual model is very similar to the visual model, since it also uses CNNs for the feature extraction process, for both video and audio. A cross-attention mechanism is then applied to these extracted features, followed by a fusion block, which is gated by the reliability of each modality at any given time. Lastly, a BiLSTM architecture is again used as the Sequence Encoder.

The experimental process of this work consists of three separate experiments. During the first, the visual model is benchmarked in a constrained environment, in terms of both the vocabulary spoken by the speaker as well as the syntax they employ. The first sub-experiment is an ablation study, which produces the final proposed method as the model with the best results. The second sub-experiment studies the performance of the proposed method more in-depth, also using an algorithm for word prediction based on the predicted phonemes, which takes advantage of the constraints present in the chosen dataset. The second experiment tested the performance of the proposed audio-visual method, in the same constrained environment, but this time including noise in both modalities. During the first sub-experiment, the results of increasing audio SNR are studied, whereas during the second sub-experiment, the model is tested during conditions where the visual noise is increased. Finally, the third experiment compares the proposed audiovisual model to four state of the art unimodal models, in an environment with noisy audio as well as corrupted video. 

\vspace*{\fill}
\end{english}
